\documentclass[11pt]{article}

\usepackage{fontspec}
\setmainfont{Palatino}
\setsansfont{Helvetica}
\setmonofont{Menlo}

\usepackage{kotex}
\setmainhangulfont{Nanum Myeongjo}

\usepackage{amsmath, amssymb}
\usepackage[margin=1in]{geometry}
\usepackage{setspace}
\onehalfspacing
\usepackage{float}
\usepackage{graphicx}
\usepackage{booktabs}
\usepackage{xcolor}
\definecolor{blue}{RGB}{0,114,178}
\usepackage{hyperref}
\hypersetup{colorlinks, citecolor=blue, linkcolor=blue, urlcolor=blue}

\begin{document}

\title{When Does Policy Content Matter? Committee Processing of\\
Redistributive and Distributive Legislation\\
in the Korean National Assembly}
\author{KNA Research Agents (AI-generated)\thanks{Experimental Output - kna-research-agents.com}}
\date{\today}
\maketitle
\thispagestyle{empty}

\begin{abstract}
\noindent Does the substantive content of legislation shape its fate in committee? Despite six decades of theorizing since \citeauthor{lowi1964}'s (\citeyear{lowi1964}) policy typology, to my knowledge no study has tested whether redistributive bills face systematically different committee processing outcomes than distributive bills at the individual bill level. I address this gap using a keyword classification of over 40,000 member-sponsored bills in the Korean National Assembly across five assemblies (17th--21st, 2004--2024). Labor bills, the paradigmatic redistributive category, receive committee decisions at substantially lower rates than small-business bills, the corresponding distributive category, after controlling for committee assignment, sponsor characteristics, and institutional access. This content-based gradient persists after controlling for committee membership, which appears to function as a separate institutional gate. Committee-restricted specifications, which reduce classification noise, amplify rather than attenuate the estimated gradient. Across five assemblies, the gradient's direction is stable from the 18th Assembly onward, but its magnitude varies with regime type, a suggestive descriptive pattern consistent with the possibility that Lowi's content friction is politically contingent.
\end{abstract}

\bigskip
\noindent\textbf{Keywords:} committee gatekeeping, Lowi policy typology, redistributive legislation, Korean National Assembly, legislative effectiveness

\bigskip
\setcounter{page}{0}
\clearpage

%% ============================================================
%% INTRODUCTION
%% ============================================================

\section{Introduction}
\label{sec:intro}

\citet{lowi1964} proposed that the substance of policy determines the structure of politics surrounding it. Redistributive policies, which transfer resources across broad social classes, generate organized opposition from identifiable losers; distributive policies, which allocate particularistic benefits without imposing visible costs, attract logrolling coalitions. This distinction implies that legislatures should process the two types of legislation differently. Redistributive bills should face more friction because they activate opposition, while distributive bills should advance more smoothly because they serve concentrated interests. Yet to my knowledge, no study has tested this prediction at the level of individual bills within the committee stage of any legislature.

The absence of bill-level evidence for Lowi's prediction is puzzling given the centrality of his typology in comparative politics. Two practical barriers may explain this lacuna. First, classifying thousands of bills by policy type requires either prohibitively expensive hand-coding or keyword-based approaches whose accuracy is difficult to verify. Second, committees handle bills from multiple policy domains, making it difficult to isolate content-based processing differences from committee-specific institutional effects. Most empirical work on committee behavior has accordingly focused on partisan \citep{cox2005}, informational \citep{krehbiel1991}, or sponsor-level \citep{volden2014} determinants of bill survival, leaving the role of substantive policy content largely unexamined.

This paper addresses both barriers using data from the Korean National Assembly (KNA), where a comprehensive digital record of legislative activity provides the raw material for systematic bill classification, and where I introduce a committee-routing validation method that independently confirms classifier accuracy. The KNA is an especially informative setting because approximately 80 percent of failed member-sponsored bills in recent assemblies died at a single chokepoint: they were placed on a standing committee agenda but never received a committee decision. Understanding what predicts which bills survive this gate is therefore central to understanding legislative outcomes in Korea's democracy.

I classify member-sponsored bills into Lowi's redistributive and distributive categories using strict keyword matching applied to bill names and propose-reason texts. Labor bills (those addressing minimum wages, labor standards, industrial accidents, and union rights) represent the redistributive category; small-business bills (those targeting SMEs, traditional markets, and micro-enterprises) represent the distributive category. To validate this classification, I check whether each bill is assigned to the committee predicted by its content. Bills classified as ``Labor'' that are routed to the Environment and Labor Committee (\textit{Hwangyeong Nodong Wiwonhoe}, 환경노동위원회) receive external confirmation of the classification. Critically, restricting the sample to these correctly routed bills amplifies the estimated content gradient, confirming that keyword misclassification attenuates rather than inflates the finding.

The analysis reveals a substantial and robust gap in committee processing rates between these two categories - a pattern I term the Lowi gradient. Labor bills receive committee decisions at markedly lower rates than small-business bills, a disparity that persists across specifications controlling for committee assignment, sponsor characteristics, and institutional access (Table~\ref{tab:main}). Within-sponsor comparisons confirm that the gradient is not attributable to unobserved legislator quality. Institutional access through committee membership is associated with a separate processing advantage that does not, however, eliminate the content-based gradient: insiders who sponsor labor bills still show substantially lower processing rates than insiders who sponsor small-business bills.

The gradient's direction is consistent from the 18th Assembly onward, but its magnitude varies across assemblies in a pattern that covaries with regime type. Under conservative presidencies, the gap between labor and small-business processing rates widens dramatically; under progressive unified government, it compresses or reverses. This suggestive descriptive pattern is consistent with conditional party government theory \citep{aldrich2001}, though the small number of regime transitions (five assemblies) limits the strength of this inference and motivates future research with more variation in political context.

The remainder of the paper proceeds as follows. Section~\ref{sec:lit} situates this study within the literatures on committee organization, Lowi's policy typology, and content-specific legislative penalties, deriving theoretical expectations. Section~\ref{sec:method} describes the KNA bill database, the keyword classification approach, and the identification strategy. Section~\ref{sec:results} presents the Lowi gradient across specifications, the insider-outsider decomposition, and the cross-assembly variation. Section~\ref{sec:discussion} discusses the findings' implications, limitations, and connections to prior work. Section~\ref{sec:conclusion} concludes.

%% ============================================================
%% LITERATURE AND THEORY
%% ============================================================

\section{Literature and Theory}
\label{sec:lit}

\subsection{Committee Organization and Gatekeeping}
\label{sec:lit-committee}

Theories of legislative organization disagree on why committees exist and whom they serve, but converge on the observation that committees exercise substantial discretion over which bills advance. Under cartel theory, the majority party uses committee assignments and agenda control to prevent bills opposed by the party median from reaching the floor \citep{cox2005}. Under the informational theory, committees serve the legislature as a whole by developing policy expertise, and gatekeeping reflects the costs of processing low-information bills \citep{krehbiel1991}. \citet{krutz2005} offers a third account rooted in bounded rationality: committees winnow bills using observable cues such as sponsor identity and timing to triage an unmanageable workload.

All three theories predict that most bills will fail at the committee stage, but they differ on which bills survive. Cartel theory predicts partisan selection; informational theory predicts expertise-driven selection; winnowing theory predicts cue-based selection. None makes a strong prediction about the \emph{substantive content} of legislation as an independent determinant of committee action. A labor bill and a small-business bill introduced by the same legislator, referred to the same committee, at the same point in the legislative calendar, should receive similar treatment under all three frameworks. If they do not, a content-based mechanism may be at work that existing theories do not fully capture.

\subsection{Lowi's Policy Typology and the Empirical Gap}
\label{sec:lit-lowi}

\citet{lowi1964} argued that the content of policy determines the political dynamics surrounding it. Redistributive policies, which reallocate resources across broad social categories (labor regulation, progressive taxation), generate organized opposition from groups who stand to lose. Distributive policies, which confer particularistic benefits without imposing visible costs (small-business subsidies, infrastructure projects), attract logrolling coalitions in which legislators exchange support. The implication for committee processing is direct: redistributive bills should face greater political friction because they activate opposition, while distributive bills should advance with less resistance.

Despite the intuitive appeal of this prediction, the empirical literature on Lowi's typology has remained largely theoretical and taxonomic. Scholars have debated the boundaries of Lowi's categories and proposed refinements, but there exists, to my knowledge, no study that classifies individual bills by Lowi's redistributive-distributive distinction and examines whether they face systematically different committee processing rates. This lacuna may stem from the practical difficulty of classifying thousands of bills by policy type and the challenge of separating content effects from committee-specific institutional dynamics within any single legislature.

\subsection{Content-Specific Legislative Penalties}
\label{sec:lit-content}

Two recent studies in the United States Congress demonstrate that the substantive content of legislation can shape its legislative fate, independent of sponsor characteristics. \citet{volden2016} find that bills addressing women's issues face a processing penalty relative to other legislation, even after controlling for sponsor attributes captured by the Legislative Effectiveness Score framework \citep{volden2014}. \citet{peay2020} extends this finding to Black lawmakers, showing that committee membership does not equalize processing outcomes for all policy types; Black legislators who serve on relevant committees still face a content-specific penalty for bills addressing racial issues. These studies establish the existence of content-based legislative friction, but neither connects the finding to Lowi's typology. The penalties documented by \citeauthor{volden2016} and by \citeauthor{peay2020} may reflect the redistributive character of the disadvantaged content categories, but the authors do not test this interpretation.

While prior work on the KNA has identified sponsor-level predictors of committee outcomes - \citet{kim2023} demonstrate that a bill sponsor's subcommittee position significantly predicts passage, and \citet{an2025} identify additional sponsor-level factors in recent assemblies - none of these studies examines policy content as an independent predictor. Similarly, research on the partisan architecture of committee assignments \citep{choi2018} and the scrutiny process for politically controversial bills \citep{seo2020, jeong2024} has illuminated institutional mechanisms of legislative processing without addressing whether the substance of legislation, independent of its political salience or sponsor characteristics, predicts committee outcomes. The present study bridges this gap by testing Lowi's content-based prediction directly.

\subsection{The Korean National Assembly as a Testing Ground}
\label{sec:lit-kna}

The KNA provides an unusually favorable setting for testing Lowi's prediction at the bill level. First, the volume of member-sponsored legislation has grown dramatically, from approximately 5,700 bills in the 17th Assembly (2004--2008) to over 25,000 in the 21st Assembly (2020--2024), creating substantial variation in processing outcomes. Second, the KNA's bill-tracking system records the complete legislative trajectory of each bill, from introduction through committee referral, committee decision, and plenary vote, enabling precise measurement of the committee processing stage where most bill attrition occurs. Third, the KNA's standing committee structure maps relatively cleanly onto policy domains, allowing an external validation of content classifications through committee-routing checks.

Korea's recent political history also provides useful variation in regime type. The five completed assemblies since 2004 span two progressive presidencies (Roh Moo-hyun, Moon Jae-in), two conservative presidencies (Lee Myung-bak, Park Geun-hye), and one mixed period (the 21st Assembly, which transitioned from Moon's unified government to Yoon Suk-yeol's divided government in 2022). \citet{aldrich2001} predict that the majority party exercises stronger legislative control when its members are ideologically homogeneous and distant from the minority party. If conservative governments are ideologically opposed to labor redistribution, the Lowi gradient should intensify under conservative regimes: committees should process redistributive legislation even less favorably when the governing coalition opposes the policy domain. This prediction connects Lowi's typology to the conditional party government literature, suggesting that the intensity of content-based friction may depend on the political alignment between the governing coalition and the redistributive content of legislation.

\subsection{Theoretical Expectations}
\label{sec:lit-expect}

Building on Lowi's typology, the content-penalty literature, and the KNA's institutional features, I derive three expectations:

\medskip
\noindent\textbf{H1.} \textit{Redistributive legislation (labor bills) faces lower committee processing rates than distributive legislation (small-business bills), controlling for committee assignment, sponsor characteristics, and institutional access.}

\medskip
\noindent\textbf{H2a.} \textit{Institutional access through committee membership is associated with a processing advantage regardless of bill content.}

\medskip
\noindent\textbf{H2b.} \textit{The committee-membership advantage does not eliminate the redistributive-distributive gradient; the content-based gap persists even among committee insiders.}

\medskip
\noindent H1 follows from Lowi's prediction that redistributive content generates organized opposition at the committee stage. H2a and H2b follow from the theoretical independence of institutional access (whether partisan, informational, or positional) and content friction. If all three expectations hold, committee processing operates through two gates: an institutional gate set by the sponsor's relationship to the committee, and a content gate set by the policy type of the legislation.

%% ============================================================
%% DATA AND METHOD
%% ============================================================

\section{Data and Method}
\label{sec:method}

\subsection{Data}
\label{sec:data}

I draw on the KNA bill database, which provides comprehensive records of all legislation introduced since the 17th Assembly (2004). The primary analysis uses member-sponsored bills (\textit{uiwon barui beomnyulan}, 의원발의 법률안) from the 20th and 21st Assemblies (2016--2024), with a cross-assembly extension covering all five completed assemblies (17th--21st, 2004--2024). I exclude government-sponsored and committee-sponsored bills, which follow different processing pathways. The 20th and 21st Assemblies together produced approximately 46,000 member-sponsored bills.

The \textbf{dependent variable} is committee decision, coded 1 if the standing committee took any action on the bill (including passage, rejection, or alternative-bill incorporation under the \textit{daean banyeong pyegi} (대안반영폐기) procedure) and 0 if the bill expired without receiving a committee decision. This binary measure captures the gatekeeping stage where the vast majority of bill attrition occurs: roughly 80 percent of failed bills in recent assemblies died after being placed on a committee agenda without ever receiving a decision.

I classify bills into Lowi's categories using a \textbf{strict keyword classifier} applied to bill names and, where available, propose-reason texts (available only for the 20th and 21st Assemblies). The redistributive category consists of labor bills, identified by keywords including \textit{choejeo imgeum} (최저임금, minimum wage), \textit{geunro gijun} (근로기준, labor standards), \textit{nodong johap} (노동조합, labor union), \textit{saneop jaehae} (산업재해, industrial accident), \textit{goyong boheom} (고용보험, employment insurance), \textit{bijeong-gyu} (비정규, non-regular work), and related terms. The distributive category consists of small-business bills, identified by keywords including \textit{sosanggo-in} (소상공인, small business owner), \textit{jungso-gieop} (중소기업, SME), \textit{jeontong sijang} (전통시장, traditional market), \textit{sangga imdae} (상가임대, commercial lease), and related terms. Two broad keywords, 근로자 (worker) and 노동자 (laborer), are excluded from the strict classifier because they appear in bills across many policy domains, reducing the classifier's committee-alignment rate.

To validate the classifier, I introduce a \textbf{committee-routing check}. If a bill classified as labor legislation is assigned to the Environment and Labor Committee (\textit{Hwangyeong Nodong Wiwonhoe}, 환경노동위원회), the classification is externally confirmed. The strict classifier achieves approximately 58 percent committee alignment for labor bills and 78 percent for welfare bills. Committee-restricted specifications, which limit the sample to correctly routed bills, produce \emph{larger} content gradients than unrestricted specifications, confirming that misclassification attenuates rather than inflates the estimated processing gap (consistent with classical measurement error in the independent variable).

Table~\ref{tab:rates} presents committee decision rates by policy category across all five completed assemblies. The strict classifier produces a Lowi sample of approximately 2,800 labor bills and 2,600 small-business bills in the 20th and 21st Assemblies.

The key control variable is \textbf{sponsor-committee match}, coded 1 if the bill's primary sponsor is identified as serving on the committee to which the bill is referred. I proxy committee membership using the modal committee of each legislator's bill referrals within each assembly term.\footnote{This proxy relies on the assumption that legislators introduce bills disproportionately to their own committee. However, legislators who pursue policy interests outside their committee assignment could be misclassified. I was unable to validate this proxy against official KNA committee rosters for the analysis period. Under classical (non-differential) measurement error, the proxy biases the committee-match coefficient toward zero without biasing the content coefficient. If measurement error is differential - for instance, if labor-bill sponsors are more likely to be misclassified - the content gradient could be biased in either direction. This limitation warrants validation against official rosters in future work.} Additional controls include arrival timing (year of introduction within the assembly term) and bill text length as a proxy for complexity.

\subsection{Identification Strategy}
\label{sec:ident}

I estimate a linear probability model of committee decision as a function of policy content, institutional access, and controls:

\begin{equation}
\text{Decision}_i = \alpha + \beta_1 \text{Labor}_i + \beta_2 \text{Match}_i + \mathbf{X}_i \boldsymbol{\gamma} + \delta_c + \epsilon_i
\label{eq:main}
\end{equation}

\noindent where $\text{Decision}_i$ is a binary indicator for whether bill $i$ received any committee decision, $\text{Labor}_i$ indicates redistributive (labor) content, $\text{Match}_i$ indicates that the sponsor serves on the receiving committee, $\mathbf{X}_i$ is a vector of controls, $\delta_c$ represents committee fixed effects, and $\epsilon_i$ is the error term. The coefficient $\beta_1$ captures the average difference in committee processing rates between redistributive and distributive legislation, conditional on institutional access and committee assignment.

Identification rests on comparing labor and small-business bills that share institutional features. Committee fixed effects absorb time-invariant differences across committees in overall processing capacity. The within-sponsor comparison, which exploits variation in the policy content of bills introduced by the same legislator, addresses the concern that certain types of legislators may sponsor both more labor bills and less viable legislation. If the same legislator's labor bills process at lower rates than the same legislator's small-business bills, the gradient is unlikely to be attributable to unobserved legislator quality.

For the cross-assembly extension (17th--21st Assemblies), I interact the labor indicator with a regime variable:

\begin{equation}
\text{Decision}_i = \alpha + \beta_1 \text{Labor}_i + \beta_2 \text{Conservative}_a + \beta_3 (\text{Labor}_i \times \text{Conservative}_a) + \beta_4 \text{Match}_i + \epsilon_i
\label{eq:regime}
\end{equation}

\noindent where $\text{Conservative}_a$ indicates that assembly $a$ fell under a conservative presidency (the 18th and 19th Assemblies, under Lee Myung-bak and Park Geun-hye respectively). With only five assemblies and a binary regime indicator, the effective sample size for the interaction is five, not the number of bills. Conventional bill-level standard errors overstate precision. I therefore report exact permutation p-values computed over all $\binom{5}{2} = 10$ possible assignments of two ``conservative'' assemblies from five, providing honest inference given the small number of clusters. The minimum achievable permutation p-value is 0.10 (1 out of 10 possible assignments), which means this test cannot reach conventional significance levels by construction. Accordingly, the regime interaction should be understood as a suggestive descriptive pattern rather than a causally identified regime effect.

%% ============================================================
%% RESULTS
%% ============================================================

\section{Results}
\label{sec:results}

\subsection{Descriptive Patterns}
\label{sec:results-desc}

Table~\ref{tab:rates} presents committee decision rates for labor and small-business bills across the five completed assemblies. In the 20th and 21st Assemblies, which constitute the primary analysis sample, labor bills received committee decisions at substantially lower rates than small-business bills. This descriptive gap is large in substantive terms: a labor bill introduced in the 21st Assembly had roughly half the probability of receiving any committee attention compared to a small-business bill.

The cross-assembly variation is striking. Under Roh Moo-hyun's progressive government in the 17th Assembly, the gradient reversed: labor bills processed at higher rates than small-business bills. From the 18th Assembly onward, the gradient consistently favors small-business bills, but the magnitude varies across assemblies (see Table~\ref{tab:rates} for rates by assembly). The single-law comparison between Labor Standards Act (\textit{Geunro Gijunbeop}, 근로기준법) amendments and SME (중소기업) bills, which eliminates keyword classification noise entirely, shows an even sharper pattern.

\begin{table}[H]
\centering
\caption{Committee Decision Rates by Policy Type and Assembly}
\label{tab:rates}
\begin{tabular}{llccccc}
\toprule
 & & \multicolumn{2}{c}{Labor (Redist.)} & \multicolumn{2}{c}{SmallBiz (Distrib.)} & \\
\cmidrule(lr){3-4} \cmidrule(lr){5-6}
Assembly & President & Rate & $N$ & Rate & $N$ & Gradient \\
\midrule
17th & Roh Moo-hyun        & 82.5\% & 114   & 57.6\% & 59    & $+24.9$ pp \\
18th & Lee Myung-bak       & 17.2\% & 198   & 48.9\% & 139   & $-31.7$ pp \\
19th & Park Geun-hye       & 12.8\% & 336   & 72.8\% & 217   & $-60.0$ pp \\
20th & Moon Jae-in         & 27.1\% & 1,355 & 45.9\% & 1,138 & $-18.8$ pp \\
21st & Moon then Yoon      & 24.9\% & 1,450 & 44.7\% & 1,469 & $-19.8$ pp \\
\midrule
\multicolumn{2}{l}{\textit{Committee-Restricted}} & & & & & \\
20th & Moon Jae-in         & 23.8\% & 798   & 47.2\% & 379   & $-23.4$ pp \\
21st & Moon then Yoon      & 18.5\% & 834   & 46.3\% & 503   & $-27.9$ pp \\
\bottomrule
\multicolumn{7}{l}{\footnotesize Note: Strict keyword classifier, excluding broad 근로자/노동자.} \\
\multicolumn{7}{l}{\footnotesize Committee-restricted: Labor bills in 환경노동위, SmallBiz in 산업/중소위.} \\
\multicolumn{7}{l}{\footnotesize All gradients significant at $p < 0.001$ by chi-squared test.} \\
\end{tabular}
\end{table}

The bottom panel of Table~\ref{tab:rates} reports committee-restricted decision rates for the 20th and 21st Assemblies. Restricting to bills assigned to the predicted receiving committee increases the gradient by four to eight percentage points relative to the unrestricted sample. This amplification pattern is consistent across both assemblies and constitutes the paper's key measurement validation: if keyword misclassification were generating the gradient, restricting to correctly classified bills should shrink it, not enlarge it.

\subsection{Main Regression Estimates}
\label{sec:results-reg}

Table~\ref{tab:main} presents linear probability model estimates of the Lowi gradient. Column~(1) reports the baseline specification on the full five-assembly sample (17th--21st), including the regime interaction and the sponsor-committee match control. The coefficient on the labor indicator captures the Lowi gradient under non-conservative regimes. The interaction term captures the additional processing gap under conservative presidencies. Column~(2) excludes the small-sample 17th Assembly to verify that the results are not driven by the reversal under Roh Moo-hyun. Column~(3) restricts to bills assigned to the expected receiving committee, reducing classification noise at the cost of reduced sample size.

\begin{table}[H]
\centering
\caption{Linear Probability Model: The Lowi Gradient in Committee Processing}
\label{tab:main}
\begin{tabular}{lccc}
\toprule
 & (1) & (2) & (3) \\
 & Full Sample & Excl.\ 17th & Cmte-Restricted \\
 & (17--21) & (18--21) & (17, 19--21) \\
\midrule
Labor (Redistributive)  & $-0.181^{***}$ & $-0.199^{***}$ & $-0.222^{***}$ \\
                         & (0.013)        & (0.013)        & (0.018)        \\[4pt]
Conservative Regime      & $0.158^{***}$  & $-$            & $0.321^{***}$  \\
                         & (0.026)        &                & (0.036)        \\[4pt]
Labor $\times$ Conserv.  & $-0.322^{***}$ & $-0.302^{***}$ & $-0.425^{***}$ \\
                         & (0.034)        & (0.034)        & (0.042)        \\[4pt]
Sponsor-Cmte Match       & $0.085^{***}$  & $-$            & $0.114^{***}$  \\
                         & (0.013)        &                & (0.015)        \\
\midrule
$N$                      & 6,475          & 6,302          & 3,361          \\
$R^2$                    & 0.070          & $-$            & $-$            \\
Sponsor Match Controlled & Yes            & Yes            & Yes            \\
Permutation $p$ (interaction) & $-$       & $-$            & 0.10           \\
\bottomrule
\multicolumn{4}{l}{\footnotesize $^{*}p<0.10$, $^{**}p<0.05$, $^{***}p<0.01$. Robust SE in parentheses.} \\
\multicolumn{4}{l}{\footnotesize Conservative = 18th \& 19th Assemblies (Lee Myung-bak, Park Geun-hye).} \\
\multicolumn{4}{l}{\footnotesize Col.\ (2) omits Conservative Regime and reports the Labor $\times$ Conserv.} \\
\multicolumn{4}{l}{\footnotesize interaction without the main regime term. ``$-$'' = not included/reported.} \\
\multicolumn{4}{l}{\footnotesize Permutation $p$: exact test over all $\binom{5}{2} = 10$ possible regime assignments.} \\
\end{tabular}
\end{table}

The results are consistent across specifications. The Lowi gradient under non-conservative regimes is associated with a processing difference of roughly 18 to 22 percentage points, with labor bills processing at substantially lower rates than small-business bills (Table~\ref{tab:main}). The regime interaction is large and negative across all three columns: the estimated gradient widens substantially under conservative presidencies. In the committee-restricted specification (Column~3), which provides the cleanest measurement, the conditional gradient under conservative regimes reaches a substantively enormous magnitude.

The exact permutation test reported in Column~(3) shows that the observed regime interaction is the most extreme of all 10 possible assignments of two ``conservative'' assemblies from five (permutation $p = 0.10$, one-sided). While this is the most favorable outcome achievable given the design, it does not reach conventional significance. The regime interaction is therefore best understood as a descriptive pattern - one that motivates targeted investigation with more regime transitions - rather than a finding with strong statistical support. Additional assemblies, or within-assembly temporal variation such as mid-term government transitions, would provide the variation needed for more rigorous inference.

\subsection{The Insider-Outsider Decomposition}
\label{sec:results-insider}

A key question is whether committee insiders - those who sponsor bills referred to their own committee - show the same content-based processing disparity as outsiders. If institutional access is associated with the disappearance of the content gradient, the Lowi mechanism may operate through access channels rather than content-based friction per se.

The data suggest a nuanced pattern. Committee insiders show no statistically significant average processing disparity for livelihood legislation as a whole (approximately two percentage points, not distinguishable from zero). Outsiders show a significant gap of approximately seven percentage points. This pattern is consistent with H2a: institutional access is associated with a processing advantage that appears to neutralize the average content friction.

However, the Lowi gradient between labor and small-business bills persists for both insiders and outsiders at nearly identical magnitudes (approximately 18 percentage points for each group). Committee membership is associated with the elimination of the average content penalty but not with the elimination of the steep redistributive-distributive gradient. A legislator who serves on a committee shows lower processing rates for labor bills than for small-business bills, even when both are referred to that legislator's own committee. This finding supports a two-gate model, consistent with both H2a and H2b: institutional access (Gate 1) is associated with processing advantages independent of policy content, while the Lowi gradient (Gate 2) is associated with differential friction based on the redistributive character of the legislation.

An \citet{oster2019} robustness test confirms the stability of the content coefficient to unobserved selection: the proportional selection parameter $\delta$ exceeds 1.9, well above the recommended threshold of 1.0, indicating that unobservable confounders would need to be nearly twice as important as the observed controls to fully explain the gradient.

\subsection{Cross-Assembly Variation}
\label{sec:results-cross}

Figure~\ref{fig:gradient} displays the Lowi gradient across all five completed assemblies, with approximate 95 percent confidence intervals computed from the proportional difference in decision rates. The regime contingency is visually apparent: the two conservative-presidency assemblies (18th and 19th) show substantially larger negative gradients than the non-conservative assemblies, with the 19th Assembly under Park Geun-hye representing the extreme case. The 17th Assembly under Roh Moo-hyun is the sole reversal, the only period in the data where labor bills were favored over small-business bills.

\begin{verbatim}
# Figure 1: Cross-Assembly Lowi Gradient
library(ggplot2)
df <- data.frame(
  assembly = c("17th\n(Roh)", "18th\n(Lee MB)",
               "19th\n(Park GH)", "20th\n(Moon)",
               "21st\n(Moon, Yoon)"),
  gradient = c(24.9, -31.7, -60.0, -18.8, -19.8),
  se = c(7.36, 5.02, 3.53, 1.91, 1.72),
  regime = c("Progressive", "Conservative",
             "Conservative", "Progressive", "Mixed")
)
df$ci_low <- df$gradient - 1.96 * df$se
df$ci_high <- df$gradient + 1.96 * df$se
df$assembly <- factor(df$assembly, levels = df$assembly)

ggplot(df, aes(x = gradient, y = assembly, color = regime)) +
  geom_pointrange(aes(xmin = ci_low, xmax = ci_high),
                  size = 0.8) +
  geom_vline(xintercept = 0, linetype = "dashed",
             color = "gray50") +
  scale_color_manual(values = c(
    "Conservative" = "#D55E00",
    "Progressive" = "#0072B2",
    "Mixed" = "#009E73")) +
  labs(x = "Lowi Gradient: Labor - SmallBiz (pp)",
       y = NULL, color = "Regime") +
  theme_bw(base_size = 11) +
  theme(legend.position = "bottom")
ggsave("fig_gradient_assembly.pdf", width = 7, height = 4.5)
\end{verbatim}

\begin{figure}[H]
\centering
\fbox{\parbox{0.85\textwidth}{[Figure generated by R script above.
Coefficient plot showing the Lowi gradient (Labor $-$ SmallBiz committee decision rate,
in percentage points) by assembly, colored by regime type.
17th (Roh): $+24.9$ pp [10.5, 39.3], Progressive.
18th (Lee MB): $-31.7$ pp [$-41.5$, $-21.9$], Conservative.
19th (Park GH): $-60.0$ pp [$-66.9$, $-53.1$], Conservative.
20th (Moon): $-18.8$ pp [$-22.5$, $-15.1$], Progressive.
21st (Moon, Yoon): $-19.8$ pp [$-23.2$, $-16.4$], Mixed.]}}
\caption{The Lowi Gradient across Five Assemblies (17th--21st), by Regime Type}
\label{fig:gradient}
\end{figure}

Two features of Figure~\ref{fig:gradient} merit emphasis. First, the gradient's \emph{direction} is consistent from the 18th Assembly onward: labor bills always show a processing disadvantage relative to small-business bills. The 17th Assembly reversal under Roh Moo-hyun is the exception, and it occurs in an environment of low bill volume (approximately 5,700 total bills, compared to over 25,000 in the 21st) and high overall passage rates (51 percent). The reversal may partly reflect capacity conditions rather than pure regime favoritism.

Second, the gradient's \emph{magnitude} covaries with regime type, a descriptive pattern worth highlighting despite the limited number of assemblies. The two conservative-presidency assemblies cluster at the extreme negative end, while the progressive and mixed assemblies show a more moderate gap. This pattern is consistent with the prediction from conditional party government theory \citep{aldrich2001} that content-specific processing disparities should intensify when the governing coalition's ideological distance from the policy domain is greatest. However, with only five assemblies and the inherent power constraint of the permutation test (Section~\ref{sec:ident}), this regime contingency should be treated as a suggestive pattern that warrants further investigation rather than an established empirical regularity.

\subsection{Within-Party Evidence against Pure Partisan Gatekeeping}
\label{sec:results-party}

An important alternative explanation is that the Lowi gradient reflects partisan gatekeeping (conservative committee chairs blocking labor bills from opposition sponsors) rather than content-based friction. If this were the sole mechanism, the gradient should disappear within each party's own bill portfolio. It does not. In the 20th and 21st Assemblies, the Lowi gradient persists within both partisan blocs: liberal-bloc legislators show a gap of approximately 14 to 17 percentage points between their own labor and small-business bills, while conservative-bloc legislators show an even larger gap. The persistence of the gradient within both parties suggests that content-based friction is associated with processing outcomes independently of the sponsor's partisan identity, although the larger gradient for conservative-bloc sponsors may reflect additional dynamics that I cannot fully decompose without committee leadership data.

%% ============================================================
%% DISCUSSION
%% ============================================================

\section{Discussion}
\label{sec:discussion}

The results are consistent with both theoretical expectations. Consistent with H1, redistributive legislation is associated with a systematic committee processing disadvantage relative to distributive legislation, a finding that holds across specifications, samples, and measurement approaches. Consistent with H2a and H2b, institutional access through committee membership is associated with a separate processing advantage that does not eliminate the content-based gradient. Together, these findings point toward a two-gate model of committee processing in which institutional position and policy content are independently associated with a bill's probability of receiving committee attention.

The connection to \citet{lowi1964} merits careful framing. Lowi predicted that redistributive policies activate organized opposition. The data are consistent with this prediction: labor bills, which propose to regulate wages, working conditions, and union rights, are associated with substantially worse committee outcomes than small-business bills, which distribute benefits to concentrated constituencies. But the cross-assembly variation reveals a qualification that Lowi did not articulate: the intensity of this gap appears contingent on the political regime. Under progressive unified government, the gap may compress or reverse; under conservative government, it widens to a point where committee processing of labor bills approaches near-total paralysis. This descriptive contingency suggests a connection between Lowi's content-based theory and the partisan legislative organization literature \citep{cox2005, aldrich2001}, though the small number of regime transitions in the data (five assemblies, with a permutation p-value floor of 0.10) means this connection should be understood as a motivation for future research rather than a demonstrated finding.

The insider-outsider decomposition speaks to a puzzle raised by \citet{peay2020}. \citeauthor{peay2020} found that committee membership does not equalize processing outcomes for identity-linked policy content in the U.S. Congress. The KNA data suggest a similar but more nuanced pattern: committee membership \emph{is} associated with the elimination of the average livelihood-bill processing gap (insiders show no statistically significant content-related disparity), yet it \emph{is not} associated with the elimination of the severe within-category Lowi gradient between labor and small-business bills. The two findings are compatible if one distinguishes between moderate and severe content friction. Institutional access may be sufficient to compensate for modest content barriers but insufficient to compensate for the organized opposition that Lowi predicted for genuinely redistributive legislation.

Several limitations warrant acknowledgment. First, the keyword classifier, though validated by the committee-routing check, achieves only approximately 58 percent committee alignment for the labor category. While the attenuation-bias pattern confirms that misclassification compresses the gradient, a hand-coded validation sample would strengthen the measurement claim. Second, the sponsor-committee match variable relies on a proxy (the modal committee of each legislator's bill referrals) rather than official committee rosters. This proxy has unknown error properties: legislators who introduce many bills to a committee they do not serve on (for example, because that committee handles their policy interest) would be misclassified. I was unable to validate this proxy against actual KNA committee membership lists, which are publicly available but not yet integrated into the analysis database. If measurement error in this variable is non-classical - for instance, if labor-bill sponsors are systematically more likely to be misclassified - the content gradient estimate could be biased. Future work should validate the proxy against official rosters for at least one assembly and report the accuracy rate. Third, the regime interaction rests on a comparison of five assemblies, providing descriptive evidence rather than a causally identified regime effect. The permutation p-value of 0.10 is the most favorable outcome achievable with this design, a fundamental constraint that additional bills within existing assemblies cannot resolve. Fourth, the analysis cannot identify the mechanism through which regime type covaries with the content gradient. Committee chair party data, which would allow a direct test of partisan gatekeeping through chair assignments, remains unavailable in the current database. The within-bloc gradient provides partial evidence against a pure chair-gatekeeping explanation but cannot fully resolve this ambiguity.

%% ============================================================
%% CONCLUSION
%% ============================================================

\section{Conclusion}
\label{sec:conclusion}

This paper presents what appears to be the first bill-level test of \citeauthor{lowi1964}'s (\citeyear{lowi1964}) prediction that redistributive legislation faces greater political friction than distributive legislation at the committee stage. Using keyword classification of over 40,000 member-sponsored bills in the Korean National Assembly across five assemblies (2004--2024), I document a Lowi gradient in which labor bills are associated with committee decision rates substantially below those of small-business bills (Table~\ref{tab:main}). This gradient persists across controls for committee assignment, sponsor characteristics, and institutional access. It holds within the bill portfolios of individual legislators and within both major partisan blocs.

The findings suggest that committee gatekeeping in the KNA may operate through two independent mechanisms: an institutional gate (committee membership) and a content gate (the redistributive-distributive gradient). These gates appear to function at different levels: institutional access is associated with the absence of moderate content friction but does not offset the severe processing disadvantage facing genuinely redistributive legislation. The content gate's magnitude, moreover, covaries with the political regime in a descriptively suggestive pattern, varying from a reversal under progressive unified government to near-total paralysis under conservative government - though this regime contingency rests on only five assembly-level observations and should be treated as a motivating pattern for future research.

These results carry implications for both legislative theory and democratic accountability. If redistributive legislation - the type of policy that directly addresses inequality - faces the steepest committee barriers, then the legislative process itself may contribute to what \citet{hacker2004} calls ``policy drift'': the failure to update redistributive policies in the face of changing economic conditions. The 31 minimum wage bills introduced in the 21st Assembly that received zero committee decisions illustrate the distributional cost of this content-specific friction.

Future work should pursue three directions. First, obtaining official committee rosters and chair party data would allow a direct test of whether the regime-correlated variation operates through partisan appointment of committee chairs or through broader changes in the political environment. Second, extending the keyword classification to other legislatures would test whether the Lowi gradient is specific to the Korean institutional context or generalizable across legislative systems with strong committee structures. Third, connecting the committee-level content gradient to policy outcomes - specifically, whether content-specific committee paralysis produces measurable policy drift - would establish the welfare consequences of the processing disparities documented here.

\bigskip
\noindent\textit{This working paper was generated by AI research agents as an
experimental output. It has not been peer-reviewed or fact-checked.
Do not cite or use in any academic, policy, or professional context.}

%% ============================================================
%% REFERENCES
%% ============================================================

\bibliographystyle{apsr}
\bibliography{references}

\end{document}
